\chapter{Method}
\label{sec:method}
As shown in Fig.~\ref{fig:architecture}, ECASP takes as input multi-view binary masks together with camera calibration \((K, T)\).
Each mask is cropped, resized, and encoded by a shared DINOv3~\cite{simeoni2025dinov3} backbone, followed by projection to the fusion dimension to obtain a per-view feature map.
A Pl\"ucker-ray positional encoding, computed from the camera calibration, is then added to each feature map.
ECASP enriches each view through epipolar cross-view attention, which constrains cross-view interactions to corresponding epipolar lines in the other views.
The enriched per-view features are combined by weighted aggregation, producing a single multi-view representation.
This representation together with the enriched feature maps are passed to a hybrid pose head, which regresses rotation and estimates translation by combining geometric triangulation with a learned residual correction.
During training, ECASP additionally uses mask-completion and allocentric-rotation objectives.

\section{Problem Setup}
\label{sec:problem_setup}
ECASP takes as input \(N\) calibrated views \(V=\{v_i\}_{i=0}^{N-1}\) of a known rigid target object with an associated CAD model.
Each view \(v_i \in V\) provides a binary instance mask \(M_{v_i} \in \{0,1\}^{H_{v_i} \times W_{v_i}}\), camera intrinsics \(K_{v_i} \in \mathbb{R}^{3\times3}\), and an extrinsic transform \(T_{v_0\rightarrow v_i} \in SE(3)\) from the reference camera frame to view \(v_i\).
We assume a standard pinhole camera model, with pixel coordinates defined relative to the top-left corner of the image canvas.
Without loss of generality, we denote the reference view by \(v_0\), so \(T_{v_0\rightarrow v_0}=I\).
The reference camera may vary across samples, as described in Sec.~\ref{sec:reference_view_protocol}.

The network outputs the object pose \((R,t)\), where \(R \in SO(3)\) and \(t \in \mathbb{R}^3\), expressed in the reference camera frame.
The pose maps object-frame coordinates to the reference camera frame: \(X_{v_0} = R X_{\mathrm{obj}} + t\).
At inference, ECASP requires only binary masks and camera calibration.

\section{Per-View Mask Encoding}
\label{sec:encoding}

\paragraph{Preprocessing.}
For each view \(v_i\), we crop the foreground region from the input mask \(M_{v_i}\), resize it while preserving the aspect ratio, and pad it to a fixed-size canvas.
Since preprocessing changes the image coordinate system, we update the camera intrinsics accordingly and define a virtual intrinsic matrix \(\tilde{K}_{v_i}\) on the preprocessed canvas.
All subsequent geometric operations, including Pl\"ucker-ray construction, fundamental-matrix computation, epipolar sampling, and DLT triangulation, use \(\tilde{K}_{v_i}\).
To preserve spatial information affected by cropping, we define for each view a crop descriptor \(c_{v_i}=[c_x^i,c_y^i,f_x^i,f_y^i,s^i]\), where \(c_x^i\) and \(c_y^i\) encode the crop-induced principal-point offset, \(f_x^i\) and \(f_y^i\) are the focal lengths on the preprocessed canvas, and \(s^i\) is the resize scale.
The descriptors from all views are concatenated and used to condition the pose head.
Full details are provided in the supplementary material.

\paragraph{Backbone encoding.}
Each preprocessed binary mask is encoded by a shared DINOv3-pretrained ConvNeXt-Tiny backbone~\cite{simeoni2025dinov3,liu2022convnext}, adapted to the binary-mask domain with LoRA~\cite{hu2022lora} and lightweight fine-tuning.
Specifically, we fine-tune the convolutional stem and LayerNorm parameters to further adapt the backbone to binary-mask inputs without destabilizing the pretrained representation.
Since the backbone expects three input channels, the binary mask is replicated across channels and normalized using the pretrained backbone statistics.
We use the third-stage feature map and project it with a \(1\times1\) convolution to the fusion dimension \(C\), obtaining \(F_{v_i}\in\mathbb{R}^{C\times H'\times W'}\).



\paragraph{Pl\"ucker-ray positional encoding.}
Before epipolar cross-view attention, Pl\"ucker-ray positional encoding~\cite{sitzmann2021lightfield,kant2024spad} is added to each feature map \(F_{v_i}\).
For each feature location, we back-project the corresponding canvas pixel using the virtual intrinsics \(\tilde K_{v_i}\) and express the resulting ray in the reference frame using calibrated extrinsics.
The ray is represented by Pl\"ucker coordinates \((d,m)\in\mathbb{R}^6\), where \(o\in\mathbb{R}^3\) is the camera center, \(d\in\mathbb{R}^3\) is the normalized ray direction, and \(m=o\times d\) is its moment.
A small MLP maps the ray representation to \(C\) channels and adds it to \(F_{v_i}\), injecting calibrated camera-ray geometry.
The exact coordinate-frame convention is provided in the supplementary material.


\section{Epipolar Cross-View Attention.}
\label{sec:fusion}

\paragraph{Overview.}
The epipolar cross-view attention module enriches the \(N\) per-view feature maps \(\{F_{v_i}\}_{i=0}^{N-1}\) by exchanging information across views, with cross-view interactions constrained by epipolar geometry.
This constrains the correspondence search without requiring explicit depth hypotheses or a cost volume.
For a query view \(v_i\), tokens from \(F_{v_i}\) attend to features sampled from the remaining feature maps \(\{F_{v_j}\}_{j\ne i}\) along the corresponding calibrated epipolar lines (Fig.~\ref{fig:epipolar_attention}).
Each view is treated as a query view, producing enriched feature maps \(\{\hat{F}_{v_i}\}_{i=0}^{N-1}\).


\paragraph{Epipolar sampling.}
For each query-source pair \((v_i,v_j)\), with \(i \neq j\), we compute the relative transform:
\begin{equation}
T_{v_i\rightarrow v_j}=T_{v_0\rightarrow v_j}T_{v_0\rightarrow v_i}^{-1}.
\end{equation}
From this transform and the virtual intrinsics \(\tilde{K}_{v_i},\tilde{K}_{v_j}\),
we form the Fundamental matrix \(\mathbf{F}_{ij}\)~\cite{hartley2004multipleview}.
For each homogeneous query location \(\mathbf{p}\) on the canvas of view \(v_i\), the corresponding epipolar line in source view \(v_j\) is \(\ell_{ij}(\mathbf{p})=\mathbf{F}_{ij}\mathbf{p}\).
We sample \(S\) points uniformly along the segment where this line intersects the source-view canvas bounds, as illustrated in Fig.~\ref{fig:epipolar_attention}.
The sampled canvas coordinates are mapped to the feature grid, and source features are obtained by bilinear interpolation of \(F_{v_j}\).
These samples provide \(S\) candidate key/value tokens per query location and source view.

\paragraph{Epipolar attention.}
For each query view \(v_i\), the \(S\) samples from each of the \(N-1\) source views are concatenated, yielding \(S(N-1)\) candidate tokens per query location.
This design is inspired by prior work on epipolar feature aggregation~\cite{he2020epipolar}.
As shown in Fig.~\ref{fig:epipolar_attention}, these tokens serve as keys and values in a multi-head cross-attention block~\cite{vaswani2017attention}, while queries are taken from \(F_{v_i}\).
Let \(f_{v_i,p}\in\mathbb{R}^{C}\) denote the query feature at location \(p\) in view \(v_i\), and let \(\mathcal{S}_{v_i,p}\in\mathbb{R}^{S(N-1)\times C}\) contain the corresponding epipolar samples from all source feature maps.
For one attention head, we compute:
\begin{equation}
\begin{gathered}
q_{v_i,p} = f_{v_i,p}W_Q, \; K_{v_i,p} = \mathcal{S}_{v_i,p}W_K, \; V_{v_i,p} = \mathcal{S}_{v_i,p}W_V, \\
\operatorname{Attn}(f_{v_i,p},\mathcal{S}_{v_i,p})
=
\operatorname{softmax}\!\left(
\frac{q_{v_i,p}K_{v_i,p}^{\top}}{\sqrt{d}}
\right)V_{v_i,p}.
\end{gathered}
\end{equation}
Here, \(W_Q\), \(W_K\), and \(W_V\) are learned query, key, and value projections, respectively, and \(d\) is the attention-head dimension.
The outputs of all attention heads are concatenated and projected back to the feature dimension.
Applying the same operation with each view as the query produces enriched feature maps \(\{\hat{F}_{v_i}\}_{i=0}^{N-1}\).

\section{Weighted Aggregation}
% \paragraph{Confidence-weighted aggregation.}
\label{sec:aggregation}
Following epipolar cross-view attention, the enriched per-view features \(\{\hat{F}_{v_i}\}_{i=0}^{N-1}\) are combined into a single multi-view representation \(F_{\mathrm{agg}}\in\mathbb{R}^{C\times H'\times W'}\) using learned weighted aggregation.
At each spatial location \((h,w)\), the model assigns one scalar weight \(w_{v_i,h,w}\) to each view.
A spatial routing head predicts \(z^{\mathrm{loc}}_{v_i,h,w}\) from the local enriched features, while a view-confidence head predicts the global logit \(z^{\mathrm{view}}_{v_i}\) for each view.
The two logits are combined and normalized across views:
\begin{equation}
w_{v_i,h,w}
=
\operatorname{softmax}_{i}\!\left(
z^{\mathrm{loc}}_{v_i,h,w}
+
z^{\mathrm{view}}_{v_i}
\right).
\end{equation}
The aggregated feature at location \((h,w)\) is then
\begin{equation}
F^{\,\prime}_{\mathrm{agg},:,h,w}=
\sum_{i=0}^{N-1}
w_{v_i,h,w}\,\hat{F}_{v_i,:,h,w},
\end{equation}
The resulting representation is then refined through a residual block:
\begin{equation}
F_{\mathrm{agg}}=
F^{\,\prime}_{\mathrm{agg}}
+
g\bigl(F^{\,\prime}_{\mathrm{agg}}\bigr),
\end{equation}
where \(g\) is a lightweight residual refinement block.
Both logit heads are zero-initialized, so aggregation starts from uniform view averaging and learns view-dependent weighting from the training objective.


\section{Hybrid Pose Head}
\label{sec:pose_head}

\paragraph{Overview.}
The hybrid pose head, illustrated in Fig.~\ref{fig:hybrid_pose_head}, estimates the object pose \((R,t)\) in the reference camera frame from the aggregated representation \(F_{\mathrm{agg}}\) and the per-view enriched features \(\{\hat{F}_{v_i}\}_{i=0}^{N-1}\).
A regression branch predicts rotation directly.
Translation uses a hybrid design: a geometric branch triangulates the CAD origin to obtain \(t_{\mathrm{tri}}\), while a learned branch predicts a per-sample gated residual correction \(r\).
The final translation combines the triangulated estimate with the gated residual.

\paragraph{Rotation branch.}
\(F_{\mathrm{agg}}\) is first processed by a rotation-specific squeeze trunk, implemented as a lightweight \(1\times1\)-convolutional bottleneck, and then modulated by an AdaLN-Zero conditioner~\cite{peebles2023dit}.
The conditioner predicts per-channel scale and shift parameters \((\gamma,\beta)\) from the concatenated crop descriptors \([c_{v_0},\ldots,c_{v_{N-1}}]\).
A learnable-query attention pool~\cite{lee2019set} summarizes the spatial tokens of the conditioned rotation feature map, producing a global feature \(f^{\mathrm{main}}\).
A separate learnable-query attention pool jointly summarizes the enriched per-view maps by treating \(\{\hat{F}_{v_i}\}_{i=0}^{N-1}\) as a single set of \(NH'W'\) tokens, producing a multi-view feature \(f^{\mathrm{view}}\in\mathbb{R}^{d}\).
This feature provides complementary rotation cues from the enriched views that may be weakened by aggregation.
The feature \([f^{\mathrm{main}},f^{\mathrm{view}},e_{\mathrm{obj}}]\), where \(e_{\mathrm{obj}}\) is a learnable object-ID embedding, is mapped by an MLP to a continuous 6D rotation representation~\cite{zhou2019rotationcontinuity}.
Gram--Schmidt orthonormalization maps the 6D representation to \(R\in\mathrm{SO}(3)\).

\paragraph{Translation residual branch.}
\(F_{\mathrm{agg}}\) is processed by a separate translation-specific squeeze trunk and AdaLN-Zero conditioner using the same concatenated crop descriptors.
The conditioned translation feature map is augmented with normalized coordinate channels, spatially reduced by convolutional layers, and flattened.
Two heads predict a lateral residual \(r_{xy}\in\mathbb{R}^{2}\) and a depth residual \(r_z\in\mathbb{R}\), which are concatenated into \(r\in\mathbb{R}^{3}\).
A learned per-sample gate scales the residual before it is added to the geometric translation estimate, as described in Sec.~\ref{sec:confidence_gated_fusion}.

\paragraph{Geometric translation via single-point triangulation.}
The geometric branch estimates translation by triangulating the CAD origin from its predicted 2D projections.
A shared keypoint head predicts a single-channel heatmap at feature resolution \(H'\times W'\) from each enriched feature map \(\hat{F}_{v_i}\).
A spatial soft-argmax~\cite{nibali2018numerical} yields a continuous feature-grid location, which is mapped to the preprocessed canvas to obtain \(\hat{p}_{v_i}\in\mathbb{R}^{2}\).
Since the CAD origin satisfies \(X_{\mathrm{obj}}=0\), its position in the reference camera frame is exactly the object translation.
Given the predicted points \(\{\hat{p}_{v_i}\}_{i=0}^{N-1}\) and the calibrated projection matrices \(P_{v_i}=\tilde{K}_{v_i}[R_{v_0\rightarrow v_i}\mid t_{v_0\rightarrow v_i}]\), we recover the 3D point by standard DLT triangulation~\cite{hartley2004multipleview}, followed by dehomogenization.
The resulting point directly gives the geometric translation estimate \(t_{\mathrm{tri}}\in\mathbb{R}^{3}\) in the reference camera frame.

\paragraph{Gated translation residual.}
\label{sec:confidence_gated_fusion}
The residual \(r\) is combined with the geometric estimate \(t_{\mathrm{tri}}\) using a learned per-sample scalar gate,
\begin{equation}
t = t_{\mathrm{tri}} + w_{\mathrm{gate}} r,
\qquad
w_{\mathrm{gate}} = \sigma\bigl(\phi(h_{\mathrm{trans}})\bigr).
\end{equation}
Here, \(h_{\mathrm{trans}}\) is the translation-trunk feature, \(\phi\) predicts one scalar logit per sample, and \(\sigma\) maps its output to \(w_{\mathrm{gate}}\in[0,1]\).
The zero-initialized residual head starts from the triangulated estimate, while the learned gate enables sample-dependent refinement through the residual correction.

\paragraph{Auxiliary heads.}
Two auxiliary heads are used only during training and discarded at inference. 
A \emph{mask-completion head}, consisting of a small transformer decoder followed by an upsampling refinement block, predicts an \(H\times W\) amodal mask from each \(\hat{F}_{v_i}\).
An \emph{allocentric rotation head}~\cite{kundu2018drcnn}, consisting of a small MLP applied to a globally pooled feature from each \(\hat{F}_{v_i}\), predicts a per-view allocentric rotation, i.e., relative to each camera's viewing direction.
Both are supervised by the auxiliary losses in Sec.~\ref{sec:loss}.

\section{Training Objective}
\label{sec:loss}

\noindent\textbf{Overview.}
ECASP is trained end-to-end with a composite objective combining the main pose losses, auxiliary losses, and a routing-entropy regularizer:
\begin{equation}
\begin{aligned}
\mathcal{L}
&= \lambda_{\mathrm{rot}}\mathcal{L}_{\mathrm{rot}}^{\mathrm{sym}}
 + \lambda_{\mathrm{trans}}\mathcal{L}_{\mathrm{trans}}
 + \lambda_{\mathrm{kp}}\mathcal{L}_{\mathrm{kp}}
 + \lambda_{\mathrm{allo}}\mathcal{L}_{\mathrm{allo}} \\
&\quad
 + \lambda_{\mathrm{mc}}\mathcal{L}_{\mathrm{mc}}
 - \lambda_{\mathrm{ent}}H(w).
\end{aligned}
\label{eq:total_loss}
\end{equation}
Here, \(H(w)\) denotes the entropy of the per-location routing distribution from Sec.~\ref{sec:fusion}.
The mask-completion loss is linearly ramped over the first \(E_{\mathrm{warm}}\) epochs to stabilize early pose learning.

\paragraph{Main pose losses.}
\label{sec:main_pose_losses}
Rotation is supervised with a symmetry-aware geodesic loss on \(\mathrm{SO}(3)\):
\begin{equation}
\mathcal{L}_{\mathrm{rot}}^{\mathrm{sym}}(R,R^{\mathrm{gt}})
=
\min_{S\in\mathcal{G}_c}
d_{\mathrm{SO}(3)}\!\left(R,R^{\mathrm{gt}}S\right),
\end{equation}
where \(S\) is an object-frame symmetry transform in the symmetry set \(\mathcal{G}_c\). Translation is supervised in the reference camera frame using a decomposed \(\ell_1\) loss over lateral and depth components:
\begin{equation}
\mathcal{L}_{\mathrm{trans}}(t,t^{\mathrm{gt}})
=
w_{xy}\left\|t_{xy}-t_{xy}^{\mathrm{gt}}\right\|_1
+
w_z\left|t_z-t_{z}^{\mathrm{gt}}\right|.
\end{equation}

\paragraph{Auxiliary losses.}
In addition to the main pose losses, ECASP uses three auxiliary losses during training.
The keypoint-origin loss \(\mathcal{L}_{\mathrm{kp}}\) supervises the per-view 2D projections \(\{\hat{p}_{v_i}\}_{i=0}^{N-1}\) of the CAD origin using a smooth-\(\ell_1\) loss.
The allocentric rotation loss \(\mathcal{L}_{\mathrm{allo}}\) supervises the allocentric head using the symmetry-aware geodesic loss from Sec.~\ref{sec:main_pose_losses}.
The mask-completion loss \(\mathcal{L}_{\mathrm{mc}}\) combines binary cross-entropy, occluded-region binary cross-entropy, and soft Dice loss~\cite{milletari2016vnet}, where the occluded-region term is evaluated on object pixels present in the full mask but absent from the visible input mask.
Full loss definitions are provided in the supplementary material.

\paragraph{Routing-entropy regularizer.}
We regularize the per-location view-weight distribution \(w\) from Sec.~\ref{sec:fusion} with an entropy \(H(w)\).
Since this term is subtracted in Eq.~\ref{eq:total_loss}, minimizing the objective discourages collapse to a single view.
Full details are provided in the supplementary material.



% \section{Implementation details}
% \label{subsec: implemntation-details}
% \textbf{Pre-processing stage}
% When using DINO models, we resized the images to 480x480 pixels. When using DINOv2 models, we resized the images to 476x476 pixels. 

% \textbf{VoteCut}
% The complete list of the utilized models appears in the supplementary materials (see  \ref{subsec:Utilized-models}). Unless specified otherwise, we utilized all aforementioned models and set $\tau^m=0.2$ and $k_{max}=3$. Further ablations related to these hyperparameters can be found in \cref{sec:ablations}. We set $\tau^{cut}=0.15$, as done by Wang \etal \cite{wang2023cut}.

% To ensure accurate pixel alignment between proposals, we employ a two-step resizing approach. First, we resize all proposals to a fixed size prior to the IOU clustering phase. Then, after the pixel voting phase, we resize the final mask to match the image shape.

% \textbf{Training details}
% All experiments were performed using the Detectron2 \cite{wu2019detectron2} platform, using a batch size of 16 and the copy-paste augmentation \cite{Dwibedi_2017_ICCV, ghiasi2021simple}. 
% Cascade Mask R-CNN \cite{cai2018cascade} detector is used for CAD.

% All reported experiments have been conducted using two NVIDIA RTX-6000 GPUS.